\chapter*{Front Matter}
\addcontentsline{toc}{chapter}{Front Matter}

\section*{Statutory Declaration}
\draftnote{Erkl\"arung \"uber die eigenst\"andige Erbringung der Arbeit --- insert TUM's
  required wording verbatim once the official template is in place.}

\section*{Abstract}
\draftnote{English abstract, $\sim$250 words. Write it \emph{last}, after
  Chapter~\ref{ch:results} is filled, and make it match the evidence rather than the
  original project plan. The four elements it must contain, in this order:
  (i) the problem --- predicting MCI-to-AD conversion from longitudinal resting-state fMRI
  under severe sample constraint;
  (ii) the method --- staged graph pretraining plus a $\Delta t$-conditioned recurrent
  trajectory classifier, with a four-arm pre-registered encoder ablation;
  (iii) the reproducibility finding --- a published MICCAI baseline is not reproducible at
  a fixed seed, with within-seed variance exceeding between-seed variance across 19 runs,
  while this pipeline re-runs byte-identically;
  (iv) the negative architectural finding --- removing the graph encoder does not degrade
  performance, extending a published large-$N$ cross-sectional result into the
  longitudinal small-$N$ regime.
  Do not lead with an accuracy figure; at $N_{\mathrm{test}} = 34$ it is the least
  informative number in the thesis.}

\section*{Kurzfassung}
\draftnote{German translation of the abstract above. Write after the English version is
  final so the two cannot diverge.}

\section*{Acknowledgements}
\draftnote{Supervisor, lab, DELCODE consortium, and funding acknowledgement. ADNI and
  OASIS-3 both require specific data-use acknowledgement wording --- copy it verbatim from
  their respective data-use agreements rather than paraphrasing.}

\listoffigures
\listoftables

\section*{Reading This Draft}

Four kinds of annotation appear throughout, boxed and numbered, marking the structure of the
argument rather than its content: \textbf{RESEARCH GAP} (an open problem in the field),
\textbf{GREY AREA} (evidence genuinely contested), \textbf{OUR EDGE} (where this work is
defensibly ahead of published practice), and \textbf{EXPOSURE} (where this work is
vulnerable). Appendix~\ref{app:register} collects all of them into a single register with
page references. Readers short of time may find that register the fastest route to the
substance of the thesis.

\draftnote{These annotations are a drafting aid for supervision. Decide before submission
  whether they remain in the final document --- keeping the EXPOSURE items visible is
  defensible and unusual; removing them is conventional. They are controlled from
  \texttt{main.tex} and can be made invisible without touching the chapters.}

\section*{List of Acronyms \& Abbreviations}
\begin{description}[leftmargin=3.5cm,style=nextline]
  \item[AD] Alzheimer's Disease
  \item[ADNI] Alzheimer's Disease Neuroimaging Initiative
  \item[AUC] Area Under the Receiver Operating Characteristic Curve
  \item[BOLD] Blood-Oxygen-Level-Dependent
  \item[BTGT] Brain-TokenGT
  \item[C-index] Concordance Index
  \item[CSF] Cerebrospinal Fluid
  \item[CV] Cross-Validation
  \item[DELCODE] DZNE Longitudinal Cognitive Impairment and Dementia Study
  \item[DMN] Default Mode Network
  \item[FC] Functional Connectivity
  \item[FD] Framewise Displacement
  \item[FDR] False Discovery Rate
  \item[GAAE] Graph Attention Autoencoder
  \item[GAE / VGAE] Graph Autoencoder / Variational Graph Autoencoder
  \item[GAT / GATv2] Graph Attention Network (version 2)
  \item[GCN] Graph Convolutional Network
  \item[GEC] Graph Encoder Classifier
  \item[GE-LSTM / GE-GRU] Graph-Embedding LSTM / GRU trajectory classifier
  \item[GEP] Graph Embedding Perceptron
  \item[GNN] Graph Neural Network
  \item[GRCU] Graph Convolutional Recurrent Unit (EvolveGCN-H)
  \item[JEPA] Joint-Embedding Predictive Architecture
  \item[MCI] Mild Cognitive Impairment
  \item[OASIS-3] Open Access Series of Imaging Studies, release 3
  \item[PET] Positron Emission Tomography
  \item[ROI] Region of Interest
  \item[rs-fMRI] Resting-State Functional Magnetic Resonance Imaging
  \item[SCD] Subjective Cognitive Decline
  \item[SD] Standard Deviation
  \item[$\Delta t$] Inter-visit time interval
\end{description}
